
\subsubsection{Selecting The Micro Case}
The first facet concerns how designers select micro cases that can support the broader story. In this process, designers need to consider the \textit{grounding} of the selected \textit{case}, its \textit{representational type} in relation to the broader population, and its \textit{correspondence} with macro \textit{data}.

\paragraph{\textbf{\textcolor{col}{Case Grounding}}}

This dimension describes how micro cases represent real-world individuals and experiences. A case may be grounded in a \textbf{real individual} (n=48), where the narrative centers on a specific existing person and their lived experience. For example, in \textit{On Upward Mobility}~\cite{williams2022upwardmobility} (Fig.~\ref{fig:example}-A), the author uses his own childhood and family migration history as the micro case. Alternatively, a case may take the form of a \textbf{curated persona} (n=7), which does not correspond to a single individual but synthesizes shared characteristics or experiences from multiple real individuals. For example, in \textit{The Middle School Years}~\cite{chang2025middle} (Fig.~\ref{fig:example}-B) student personas are constructed from survey results of 40 million students to represent common experiences across age groups.

\paragraph{\textbf{\textcolor{col}{Representational Type}}}

This dimension describes how the situation represented by a micro case relates to the broader population. A \textbf{typical case} (n=38) aligns with common situations or dominant trends, allowing broader patterns to be illustrated through a concrete experience. For example, in \textit{Bussed Out}~\cite{gee2017bussed} (Fig.~\ref{fig:example}-C), Quinn Raber’s one-way bus trip from San Francisco to Indianapolis illustrates how U.S. cities relocate people experiencing homelessness. A \textbf{subgroup case} (n=6) corresponds to a specific group within the population and highlights how that group differs from the whole. For example, in \textit{Abortion Mazes}~\cite{diehm2024abortion} (Fig.~\ref{fig:example}-D), six individual cases represent abortion seekers from different states, showing how policy environments create different challenges. An \textbf{exceptional case} (n=5) substantially deviates from common patterns or dominant trends. For example, in \textit{Women in the House of Commons}~\cite{dsouza2018women} (Fig.~\ref{fig:example}-E), Constance Markievicz was the first elected female member of parliament but never took office due to the abstention policy, making her trajectory distinct from most women who entered Parliament. Finally, an \textbf{any case} (n=16) strategy does not rely on predefined selection criteria and allows any individual case to enter the micro narrative. For example, in \textit{Tied in Knots}~\cite{tiedinknots} (Fig.~\ref{fig:example}-F) each anonymous sexual harassment testimony is represented as a “knot,” and readers can select any knot to explore the corresponding experience.

\paragraph{\textbf{\textcolor{col}{Data Correspondence}}}

This dimension describes how micro cases correspond to macro datasets. A \textbf{record linked} (n=26) case is directly drawn from the macro dataset and linked to a specific record within it. For example, in \textit{This Is a Teenager}~\cite{chang2024teenager} (Fig.~\ref{fig:example}-G), Alex is an actual participant in the National Longitudinal Survey of Youth, and the article uses his survey records from adolescence to adulthood. An \textbf{attribute matched} (n=22) case does not belong to the dataset itself, but the individual’s attributes match a category represented in the data. For example, in \textit{How Uber and Lyft Used a Pay Loophole to Deny NYC Drivers Millions in Pay}~\cite{lung2024lockouts} (Fig.~\ref{fig:example}-H), Mohamed, a long-term Uber and Lyft driver in New York, matches the population described by the macro data in income and employment status. A \textbf{contextually related} (n=10) case does not correspond to a specific record or subgroup, but its experience occurs within the same broader phenomenon or context represented by the macro data. For example, in \textit{Small Boat in a Great Tide}~\cite{modernchronicle} (Fig.~\ref{fig:example}-I), Zhao Dachun’s life history and macro-level data on China’s healthcare and economic development jointly portray modern China’s social transformations through a shared historical context.

\subsubsection{Constructing The Micro Narrative}
The second facet concerns how designers develop the micro narrative itself from the selected case. In this process, designers need to consider which elements to select and organize as \textit{narrative content}, and how to realize them through \textit{visual representation}.

\paragraph{\textbf{\textcolor{des}{Narrative Content}}}

This dimension describes the types of case materials selected to construct the micro narrative. A micro narrative may focus on \textbf{background attributes} (n=30), such as age, occupation, family situation, or living environment, which provide contextual information. For example, \textit{An Incalculable Loss}~\cite{nyt2020incalculable} (Fig.~\ref{fig:example}-J) presents the names, ages, occupations, and family roles of deceased individuals, making the people behind mortality statistics more recognizable. It may also focus on \textbf{events and actions} (n=29), presenting specific experiences and behaviors. For example, \textit{Visualizing the Storm}~\cite{visualizingthestorm} (Fig.~\ref{fig:example}-K) depicts PinkGirl’s experience of online violence through events including the circulation of her photo, online attacks, public responses, and efforts to defend her rights. A micro narrative may further incorporate \textbf{personal voice} (n=8) by preserving an individual’s perspectives, attitudes, and emotions. For example, \textit{30 Years of American Anxieties}~\cite{blinderman2018anxieties} (Fig.~\ref{fig:example}-L) presents anonymous readers’ accounts and emotional expressions concerning family, relationships, and life difficulties. Finally, it may highlight \textbf{changing states} (n=6), showing how a person’s conditions, circumstances, or attributes change across stages or contexts. For example, \textit{This Is a Teenager}~\cite{chang2024teenager} (Fig.~\ref{fig:example}-G) follows Alex from adolescence to adulthood, presenting changes in his family environment, education, occupation, and life circumstances.

\paragraph{\textbf{\textcolor{des}{Visual Representation}}}

This dimension describes how the selected case materials are visually represented. A micro narrative may use \textbf{original material} (n=22), preserving records such as photographs, direct quotations, and videos. For example, \textit{Safaa’s Fatal Journey}~\cite{spencer2015safaa} (Fig.~\ref{fig:example}-M) uses photographs of Safaa and her family, together with interview excerpts, to present their experiences during their escape from Syria. It may also transform individuals and their experiences into an \textbf{illustrated character} (n=22), using redrawn figures and simplified or dramatized representations. For example, \textit{A Love Story}~\cite{chang2026love} (Fig.~\ref{fig:example}-N) depicts Henry and Mary as illustrated characters and condenses their journey from meeting and marriage to maintaining their relationship during the pandemic. An \textbf{icon} (n=15) provides a more generalized representation through consistent visual symbols. For example, \textit{Happy Map}~\cite{happymap2026} (Fig.~\ref{fig:example}-O) represents respondents as cartoon-like figures distinguished by hair color, clothing, and posture. Finally, a \textbf{data-encoded figure} (n=9) incorporates individual attributes or experiences into the visual form itself, allowing the representation to convey both personal information and data. For example, in \textit{Tied in Knots}~\cite{tiedinknots} (Fig.~\ref{fig:example}-F), each testimony is encoded as a “knot” with a distinct form reflecting the victim’s identity, type of experience, and personal circumstances.

\subsubsection{Connecting Micro And Macro}
The third facet concerns how designers connect micro narratives with macro data to establish relationships between the two scales. In this process, designers need to consider how narrative connections are established through \textit{micro-macro organization} and \textit{narrative position}, how visual connections are created through \textit{spatial relationship} and \textit{visual linkage}, and how interaction supports \textit{transition} between scales.

\paragraph{\textbf{\textcolor{imp}{Micro-Macro Organization}}}

This dimension describes how micro and macro narratives are organized within a data story. A \textbf{macro-led} (n=16) organization uses the macro narrative as the main storyline, with micro narratives explaining or substantiating the broader phenomenon. For example, in \textit{South Korea Halloween Disaster}~\cite{smith2022seoul} (Fig.~\ref{fig:example}-P), the Itaewon disaster timeline structures the story, while witnesses’ experiences, such as Linda’s account, show how people attempted to escape danger. A \textbf{micro-led} (n=26) organization follows the micro narrative, with macro information providing broader context and meaning. In \textit{On Upward Mobility}~\cite{williams2022upwardmobility} (Fig.~\ref{fig:example}-A), the author’s family migration experience anchors the story, while community-level social mobility data reveals the broader context behind it. A \textbf{micro-aggregated} (n=13) organization presents a broader phenomenon by aggregating, comparing, or synthesizing multiple micro narratives. In \textit{Happy Map}~\cite{happymap2026} (Fig.~\ref{fig:example}-O), more than six hundred individual happiness moments are grouped into 13 themed islands, collectively revealing shared patterns in experiences of happiness.

\paragraph{\textbf{\textcolor{imp}{Narrative Position}}}

This dimension describes the role of micro narratives in the progression of the broader narrative. As an \textbf{entry} (n=31), a micro narrative establishes the issue through a specific personal experience and introduces subsequent macro-level discussion. For example, in \textit{A Love Story}~\cite{chang2026love} (Fig.~\ref{fig:example}-N), Henry and Mary’s experience of forming a relationship leads into broader changes in American romantic relationships surrounding the pandemic. As an \textbf{elaboration} (n=46), a micro narrative supplements macro content by revealing the processes or impacts behind broader patterns. In \textit{Women in the House of Commons}~\cite{dsouza2018women} (Fig.~\ref{fig:example}-E), female parliamentarians’ experiences illustrate the process of women’s participation in British politics. As a \textbf{reflection} (n=6), a micro narrative returns to individual consequences or meanings after macro-level discussion, encouraging readers to reconsider what the data represents. In \textit{This Is a Teenager}~\cite{chang2024teenager} (Fig.~\ref{fig:example}-G), the visualization returns from broader developmental trends to Alex’s adult outcomes, showing how childhood experiences shaped his later development.

\paragraph{\textbf{\textcolor{imp}{Spatial Relationship}}}

This dimension describes how micro and macro information is organized within visual space. In a \textbf{separated} (n=34) arrangement, the two scales occupy independent areas, requiring readers to move between them to establish a connection. In \textit{South Korea Halloween Disaster}~\cite{smith2022seoul} (Fig.~\ref{fig:example}-P), the macro-level event visualization and witnesses’ experiences appear in different sections of the story. In an \textbf{embedded} (n=7) arrangement, micro narratives become local elements within a macro-level visualization. In \textit{Happy Map}~\cite{happymap2026} (Fig.~\ref{fig:example}-O), individual happiness experiences appear as human nodes within the overall distribution. A \textbf{juxtaposed} (n=19) arrangement presents information from both scales side by side, facilitating comparison. In \textit{On Upward Mobility}~\cite{williams2022upwardmobility} (Fig.~\ref{fig:example}-A), the family migration route appears alongside maps of the corresponding communities, presenting personal experience and community-level information together. A \textbf{layered} (n=5) arrangement represents different scales through visual or interaction layers within the same space. In \textit{Bussed Out}~\cite{gee2017bussed} (Fig.~\ref{fig:example}-C), macro-level migration flows are overlaid on a map that first presents an individual’s migration journey.

\paragraph{\textbf{\textcolor{imp}{Visual Linkage}}}

This dimension describes how visual cues establish recognizable correspondence between micro and macro. \textbf{Explicit annotation} (n=36) uses text, labels, or annotations to state their relationship directly. In \textit{Safaa’s Fatal Journey}~\cite{spencer2015safaa} (Fig.~\ref{fig:example}-M), annotations identify Safaa’s family as refugees fleeing Syria. \textbf{Shared encoding} (n=26) reuses colors, shapes, or icons to indicate correspondence across scales. In \textit{Herd Immunity Simulation}~\cite{stevens2020herd} (Fig.~\ref{fig:example}-Q), three individuals with different health conditions use distinct visual encodings that reappear for corresponding groups in the population view. \textbf{Visual highlighting} (n=12) emphasizes the corresponding micro case within the macro view, helping readers locate the individual in the broader dataset. In \textit{This Is a Teenager}~\cite{chang2024teenager} (Fig.~\ref{fig:example}-G), Alex’s data point is highlighted in white among numerous adolescent data points. \textbf{Connecting marks} (n=3), such as lines, arrows, or paths, directly link information across scales. In \textit{Women in the House of Commons}~\cite{dsouza2018women} (Fig.~\ref{fig:example}-E), lines connect individual parliamentarians’ stories to their corresponding positions on the timeline of women’s parliamentary participation.

\paragraph{\textbf{\textcolor{imp}{Transition}}}

This dimension describes how readers move between micro and macro scales. \textbf{Guided progression} (n=47) uses scrolling, playback, or similar interactions to introduce information across scales in a predefined sequence. In \textit{On Upward Mobility}~\cite{williams2022upwardmobility} (Fig.~\ref{fig:example}-A), continuous scrolling guides readers between the author’s migration experiences and community-level social mobility data. \textbf{Direct selection} (n=16) allows readers to access micro narratives from a macro view. In \textit{The Academia Is Tied in Knots}~\cite{tiedinknots} (Fig.~\ref{fig:example}-F), readers select individual “knots” from the overall academic field distribution to explore the corresponding experiences and comments. \textbf{Filtered exploration} (n=4) allows readers to specify attributes or categories and examine cases that meet those conditions. In \textit{U.S. Gun Deaths}~\cite{periscopic2013gundeaths} (Fig.~\ref{fig:example}-R), readers filter records by age, gender, location, and manner of death to explore corresponding individual stories. In addition, \textbf{hierarchical navigation} (n=8) enables readers to move from overview-level to fine-grained information through zooming or expansion. In \textit{What Vets Are Experiencing}~\cite{periscopicwaitwecarry} (Fig.~\ref{fig:example}-S), the national map first shows state-level counts of veterans waiting for healthcare, after which readers can select a state to explore individual waiting experiences.
